% answers_short.tex —— 2012 年·安徽卷·理科 · 简短答案速查
% 由 tools/make_short_answers.py 自 answers.tex 的 \gksolhead/\gkanswer 提取，
% 并与 2012/2012_数学真题_安徽理_答案要点.md 逐题比对后生成，勿手改（改源 MD 后重跑）。

\begin{gknotice}
\textbf{使用说明}：本节只给最终答案便于快速核对，完整推导见解析版。答案为编写者独立推导，\textbf{非官方标准答案}。
\end{gknotice}

\gkpart{选择题}
\begin{center}\begin{tabular}{ccccccccccc}
\toprule
题号 & 1 & 2 & 3 & 4 & 5 & 6 & 7 & 8 & 9 & 10 \\
\midrule
答案 & D & C & B & B & C & A & D & A & C & D \\
\bottomrule
\end{tabular}\end{center}

\gkpart{填空题}
\begin{itemize}[leftmargin=2.6em]
  \item[\textbf{11.}] $[-3,0]$。
  \item[\textbf{12.}] $92$。
  \item[\textbf{13.}] $\sqrt3$。
  \item[\textbf{14.}] $-\dfrac98$。
  \item[\textbf{15.}] ①②③。
\end{itemize}

\gkpart{解答题}
\begin{itemize}[leftmargin=2.6em]
  \item[\textbf{16.}] (1) $T=\pi$；(2) $g(x)=\begin{cases}-\dfrac14+\dfrac14\sin2x,&x\in\left[-\pi,-\dfrac\pi2\right],\\[6pt]-\dfrac14-\dfrac14\sin2x,&x\in\left(-\dfrac\pi2,0\right].\end{cases}$
  \item[\textbf{17.}] (1) $\dfrac{n(n+1)}{(n+m)(n+m+2)}$；(2) $P(X=n+2)=\dfrac14$，$P(X=n+1)=\dfrac14+\dfrac{n}{2(2n+1)}$，$P(X=n)=\dfrac{n+1}{2(2n+1)}$；$EX=n+\dfrac34+\dfrac{n}{4n+2}$
  \item[\textbf{18.}] (1) 证明见解析；(2) $5$；(3) $-\dfrac{\sqrt5}5$
  \item[\textbf{19.}] (1) 当 $0<a\leqslant1$ 时为 $2+b$；当 $a>1$ 时为 $a+\dfrac1a+b$；(2) $a=\dfrac2{\mathrm e^2}$，$b=\dfrac12$
  \item[\textbf{20.}] (1) $\dfrac{x^2}4+\dfrac{y^2}3=1$；(2) 证明见解析
  \item[\textbf{21.}] (1) 证明见解析；(2) $0<c\leqslant\dfrac14$
\end{itemize}
